\subsection{一个变体}\label{sec:11g}
现在略微修改前述构造，将条带 $Z$ 换成上半平面 $H$。以 $\mathcal P^-\operatorname{Gr}(V)$ 表示满足以下两个性质的路径 $\lambda:[0;1]\to\operatorname{Gr}(V)$ 所成的空间：
\begin{itemize}
\item $\lambda(0)$ 与 $\lambda(1)$ 横截相交。
\item 路径对 $(\lambda,\allowbreak \lambda(1))$ 在 $s=1$ 处的穿越型负定，其中第二个分量表示取值为 $\lambda(1)$ 的常值路径。
\end{itemize}
该空间上有自然的局部常值函数 $I_H$ 和线丛 $\delta_H$。为归结到先前的框架，注意上述第二个条件保证存在 $s'<1$，使对所有 $s\in[s';1)$ 都有 $\lambda(s)\pitchfork\lambda(1)$；特别地，将 $(\lambda,\allowbreak \lambda(1))$ 限制到 $[0;s']$，便得到 $\mathcal P^2\operatorname{Gr}(V)$ 的一个元素。可以选取光滑依赖于 $\lambda$ 的 $s'$，于是限制给出映射 $\mathcal P^-\operatorname{Gr}(V)\to\mathcal P^2\operatorname{Gr}(V)$。定义 $I_H,\allowbreak \delta_H$ 为 $I_Z,\allowbreak \delta_Z$ 沿此映射的拉回。更明确地，
\SourcePage{170}
用 $q_\lambda$ 简记 $q_{\lambda,\allowbreak \lambda(1)}$，则对一般位置的路径 $\lambda\in\mathcal P^-\operatorname{Gr}(V)$ 有
\begin{equation}\label{eq:11.20}
\begin{aligned}
I_H(\lambda)&=\sum_{s<1}\operatorname{sign}(q_\lambda(s)),\\
(\delta_H)_\lambda&=\bigotimes_{s<1}(\lambda(s)\cap\lambda(1))^{\operatorname{sign}(q_\lambda(s))}.
\end{aligned}
\end{equation}
在指标理论方面，仍选择前述函数 $\psi$，考虑 $E=H\times V$ 上的 $\bar\partial$ 算子，边界条件为 $F_\lambda\subset\partial H\times V=\R\times V$，$(F_\lambda)_s=\lambda(\psi(s))$。由于 $\lambda$ 的端点横截，它是非退化算子。将其记为 $D_{H,\allowbreak \lambda}$，则有：
\begin{lemma}\label{st:11.11}
$I_H(\lambda)=\operatorname{index}(D_{H,\allowbreak \lambda})$；且 $(\delta_H)_\lambda$ 与行列式线 $\det(D_{H,\allowbreak \lambda})$ 典范同构。
\end{lemma}
\begin{proof}
先看一个简单例子，即形如
\begin{equation}\label{eq:11.21}
\lambda(s)=\exp(\alpha sI_V)\lambda(0)
\end{equation}
的路径，其中 $\lambda(0)$ 为 Lagrangian 子空间，$\alpha\in(-\pi;0)$。此时没有计入\eref{11.20}的交点，所以引理断言 $D_{H,\allowbreak \lambda}$ 的指标为零，并有选定同构 $\det(D_{H,\allowbreak \lambda})\cong\R$。首先，指标与 $\alpha$ 的选择无关，故可假设 $\alpha=-\pi/2$。在此情形，可将两个 $D_{H,\allowbreak \lambda}$ 粘合成圆盘上的算子 $D_{D,\allowbreak \rho}$，其相应 Lagrangian 子空间环路满足 $\mu(\rho)=-n$；于是\eref{11.9}与\stref{引理}{11.7}给出 $2\operatorname{index}(D_{H,\allowbreak \lambda})=\operatorname{index}(D_{D,\allowbreak \rho})=0$。其次，要说明 $D_{H,\allowbreak \lambda}$ 为单射。取直和分量后，可归结到最小秩 $n=1$，此时适用\stref{引理}{11.5}；用该引理的记号，我们的路径 $\lambda$ 可补成满足 $\mu(\rho)=0$ 的环路 $\rho$，而曲面 $H$ 满足 $|\Sigma^-|=1$。显然，这一情形特别简单，也很容易找到其他同样可行的方法。

对一般情形，$\mathcal P^-\operatorname{Gr}(V)$ 中任意路径都可在不改变 $I_H$ 或 $\delta_H$ 的情况下，变形成对某个 $s'<1$ 在 $[s';1]$ 上满足\eref{11.21}的路径。因此只需考虑已有这一性质的路径。将 $(\lambda,\allowbreak \lambda(1))$ 限制在 $[0;s']$ 上，得到 $\mathcal P^2\operatorname{Gr}(V)$ 的元素，因而得到 $Z$ 上的 $\bar\partial$ 算子；其指标与行列式线由\stref{引理}{11.9}给出，因此与\eref{11.20}一致。另一方面，对应于 $\lambda|[s';1]$ 的 $H$ 上算子属于刚才讨论的情形。将 $H$ 的端与 $Z$ 的出向端连接，便可粘合这两部分，所得算子在相差无关紧要的变形的意义下就是 $D_{H,\allowbreak \lambda}$。由\eref{11.9}即得结论。
\end{proof}

\subsection{抽象膜结构}\label{sec:11h}
\SourcePage{171}
这一对类 $(\mu,\allowbreak w_2)$ 由一个到 Eilenberg–MacLane 空间乘积的映射表示：
\begin{equation}\label{eq:11.22}
\operatorname{Gr}(V)\longrightarrow K(\Z,1)\times K(\Z/2,2).
\end{equation}
拉回两个万有纤维化的乘积，得到纤维化 $\operatorname{Gr}(V)^\#\to\operatorname{Gr}(V)$，其纤维为 $\Z\times\RP^\infty$。当 $n\geq3$ 时，$\operatorname{Gr}(V)^\#$ 可与 Postnikov 分解意义下 $\operatorname{Gr}(V)$ 的 $2$ 连通“覆盖”识别。\znote{163}{抽象膜的思想是把 Lagrangian 子空间连同其两条“障碍类”一起提升。$K(\Z,1)$ 分类整系数的一维上同调类，在这里就是 Maslov 类 $\mu$；$K(\Z/2,2)$ 分类二阶的 $\Z/2$ 类，在这里是第二 Stiefel–Whitney 类 $w_2$。提升存在当且仅当这两条障碍同时消失，其纤维因而是 $\Z\times\RP^\infty$：$\Z$ 记录 Maslov 指标的选择，$\RP^\infty=K(\Z/2,1)$ 记录 $w_2$ 的抵消方式。$n\geq3$ 时 $\operatorname{Gr}(V)^\#$ 是 $2$ 连通的，故提升在低维上唯一。nLab 分别给出了这两类障碍的上同调解释。[来源：nLab "Maslov index"，https:/\allowbreak /\allowbreak ncatlab.org/\allowbreak nlab/\allowbreak show/\allowbreak Maslov+\allowbreak index（2026-11-15 访问）；nLab "twisted spin structure"，https:/\allowbreak /\allowbreak ncatlab.org/\allowbreak nlab/\allowbreak show/\allowbreak twisted+\allowbreak spin+\allowbreak structure（2026-11-15 访问）]}当 $n\leq2$ 时，它仍连通，并分别满足
\begin{equation}\label{eq:11.23}
\begin{cases}
\pi_1(\operatorname{Gr}(V)^\#)\cong\Z/2,\quad\pi_2(\operatorname{Gr}(V)^\#)=0,&n=1,\\
\pi_1(\operatorname{Gr}(V)^\#)=0,\quad\pi_2(\operatorname{Gr}(V)^\#)\cong\Z,&n=2.
\end{cases}
\end{equation}
\begin{definition*}
称 $\operatorname{Gr}(V)^\#$ 的点为\emph{抽象线性 Lagrangian 膜}，记为 $\Lambda^\#$，其底层 Lagrangian 子空间记为 $\Lambda$。
\end{definition*}
\begin{lemma}\label{st:11.12}
取环路 $\rho^\#\in\mathcal L\operatorname{Gr}(V)^\#$，令 $D_{D,\allowbreak \rho}$ 为圆盘 $D$ 上对应于底层环路 $\rho$ 的 $\bar\partial$ 算子。则 $\operatorname{index}(D_{D,\allowbreak \rho})=n$；此外，$\rho^\#$ 确定一个同构
\[
\det(D_{D,\rho})\cong\lambda^{\mathrm{top}}(\rho(0)).
\]
\end{lemma}
\begin{proof}
能够提升到 $\operatorname{Gr}(V)^\#$ 的环路必满足 $\mu(\rho)=0$，所以指标结论是\stref{引理}{11.7}的特例。映射 $\mathcal L\operatorname{Gr}(V)^\#\to\mathcal L_0\operatorname{Gr}(V)$ 的纤维同伦等价于 $\Z\times\mathcal L\RP^\infty$。忘掉第一个因子，并只考虑第二个因子的 $\pi_0$，便得到一个诱导的二重覆盖，其上可关联实线丛 $\delta_D\to\mathcal L_0\operatorname{Gr}(V)$。由构造，它由 $T(w_2)\in H^1(\mathcal L_0\operatorname{Gr}(V);\Z/2)$ 分类，因此由\stref{引理}{11.7}，它同构于纤维为
\begin{equation}\label{eq:11.24}
\det(D_{D,\rho})^\vee\otimes\lambda^{\mathrm{top}}(\rho(0))
\end{equation}
的自然线丛。为固定两者间的一个同构，只需考察 $\mathcal L_0\operatorname{Gr}(V)$ 的任意一个连通子集，例如常值环路所成的子集。回顾\stref{引理}{11.7}的证明：当 $\rho$ 为常值时，$D_{D,\allowbreak \rho}$ 为满射，核同构于 $\rho(0)$，故 $\det(D_{D,\allowbreak \rho})^\vee\otimes\lambda^{\mathrm{top}}(\rho(0))$ 有自然选定的定向。另一方面，$\mathcal L\operatorname{Gr}(V)^\#$ 在 $\rho$ 上的纤维可与 $\Z\times\mathcal L\RP^\infty$ 识别，该识别在相差同伦及 $\Z$ 因子中的平移的意义下唯一；所以选择 $\pi_0(\mathcal L\RP^\infty)\cong\pi_1(\RP^\infty)$ 的平凡元素，便赋予 $(\delta_D)_\rho$ 一个定向。结合两者，得到所需的线丛同构。最后，依定义，$\delta_D$ 拉回到 $\mathcal L\operatorname{Gr}(V)^\#$ 后典范平凡，从而给出\eref{11.24}拉回后的平凡化。
\end{proof}
\SourcePage{172}
对于底层子空间 $(\Lambda_0,\allowbreak \Lambda_1)$ 横截相交的任意一对抽象线性 Lagrangian 膜 $(\Lambda_0^\#,\allowbreak \Lambda_1^\#)$，可以如下赋予一个整数——\emph{绝对指标} $i(\Lambda_0^\#,\allowbreak \Lambda_1^\#)$——以及一个一维实向量空间——\emph{定向空间} $o(\Lambda_0^\#,\allowbreak \Lambda_1^\#)$。选择 $\operatorname{Gr}(V)^\#$ 中连接两点的路径 $\lambda^\#$，使其投影 $\lambda$ 属于 $\mathcal P^-\operatorname{Gr}(V)$，并令
\begin{equation}\label{eq:11.25}
i(\Lambda_0^\#,\Lambda_1^\#)=I_H(\lambda),\qquad o(\Lambda_0^\#,\Lambda_1^\#)=(\delta_H)_\lambda.
\end{equation}
当然，必须证明它们与 $\lambda^\#$ 的选择无关；对 $o$ 而言，是在典范同构意义下无关。以 $\Omega^-(\operatorname{Gr}(V),\allowbreak \Lambda_0,\allowbreak \Lambda_1)$ 表示 $\mathcal P^-\operatorname{Gr}(V)$ 中从 $\Lambda_0$ 到 $\Lambda_1$ 的路径所成的子空间，以 $\Omega^-(\operatorname{Gr}(V)^\#,\allowbreak \Lambda_0^\#,\allowbreak \Lambda_1^\#)$ 表示对应的 $\operatorname{Gr}(V)^\#$ 中的路径空间。它们分别同伦等价于通常的环路空间 $\Omega\operatorname{Gr}(V)$ 与 $\Omega\operatorname{Gr}(V)^\#$。当 $n\geq3$ 时，已知 $\Omega\operatorname{Gr}(V)^\#$ 单连通，立即得到所需唯一性。当 $n=2$ 时，$w_2$ 在 $\pi_2(\operatorname{Gr}(V))$ 上仍非平凡，这意味着映射 $\Z=\pi_2(\operatorname{Gr}(V)^\#)\to\pi_2(\operatorname{Gr}(V))=\Z$ 为乘以 $2$。因此，线丛 $\delta_H$ 拉回到 $\Omega^-(\operatorname{Gr}(V)^\#,\allowbreak \Lambda_0^\#,\allowbreak \Lambda_1^\#)$ 后平凡，同样得到结论。最后，当 $n=1$ 时，由\eref{11.23}有两个不同的同伦类；二者给出相同的 $i$，而对于定向空间，我们直接规定：改变 $\lambda^\#$ 的同伦类，会给出 $(\delta_H)_\lambda$ 的一个反定向自同构。一种等价而或许更系统的做法，是通过稳定化处理这些低维情形。\znote{164}{定向空间 $o(\Lambda_0^\#,\Lambda_1^\#)$ 是一个一维实向量空间，按\stref{引理}{11.11}它典范同构于路径上某个 $\bar\partial$ 算子的行列式线。给一维实空间选定一个生成元就是给它定向，所以“定向空间”这个名字恰是行列式线加上定向数据的合称：一旦为 $o$ 选定生成元（等价地，为相应的行列式线选定定向），两个膜之间的定向数据就确定了。Floer 理论中的定向问题正是把诸如此类的线丛的定向相容地选出来；nLab 关于行列式线丛的条目讨论了这类线的整体性质。[来源：nLab "determinant line bundle"，https:/\allowbreak /\allowbreak ncatlab.org/\allowbreak nlab/\allowbreak show/\allowbreak determinant+\allowbreak line+\allowbreak bundle（2026-11-15 访问）]}

设 $S$ 为带条带状端的 Riemann 曲面，配有平凡辛向量丛 $E=S\times V$，以及由映射 $\nu:\partial S\to\operatorname{Gr}(V)$ 给出的 Lagrangian 子丛 $F_\nu\subset E|\partial S$；$\nu$ 在条带状端上局部常值。后一条件意味着，对每个无穷远点 $\zeta$，有一对 Lagrangian 子空间 $(\Lambda_{\zeta,\allowbreak 0},\allowbreak \Lambda_{\zeta,\allowbreak 1})$，使 $\nu(\epsilon_\zeta(s,\allowbreak k))=\Lambda_{\zeta,\allowbreak k}$，其中 $(s,\allowbreak k)\in\R^\pm\times\{0;1\}$。假设每一对均满足 $\Lambda_{\zeta,\allowbreak 0}\pitchfork\Lambda_{\zeta,\allowbreak 1}$，从而相应的 $\bar\partial$ 算子 $D_{S,\allowbreak \nu}$ 非退化。此外，假设给定 $\nu$ 到 $\operatorname{Gr}(V)^\#$ 的提升 $\nu^\#$，它在各端上也局部常值，因而为这些 Lagrangian 子空间提供选定原像 $\Lambda_{\zeta,\allowbreak k}^\#$。
\begin{proposition}\label{st:11.13}
以 $\hat C_1,\allowbreak \ldots,\allowbreak \hat C_k$ 表示紧化 $\hat S$ 的边界分支，令 $e_j=|\Sigma^-\cap\hat C_j|$ 为每个分支上的入向无穷远点数。对每个 $j$ 固定一个非无穷远点 $z_j\in\hat C_j$。则
\begin{equation}\label{eq:11.26}
\begin{aligned}
\operatorname{index}(D_{S,\nu})&=n(\chi(S)-|\Sigma^-|)+\sum_{\zeta\in\Sigma^\pm}\mp i(\Lambda_{\zeta,0}^\#,\Lambda_{\zeta,1}^\#),\\
\det(D_{S,\nu})&\cong\bigotimes_j\left(\lambda^{\mathrm{top}}(\nu_{z_j})^{\otimes(1-e_j)}
\otimes\bigotimes_{\zeta\in\Sigma^\pm\cap\hat C_j}o(\Lambda_{\zeta,0}^\#,\Lambda_{\zeta,1}^\#)^{\mp1}\right).
\end{aligned}
\end{equation}
\end{proposition}
\SourcePage{173}
\begin{proof}
已证明两个特殊情形。$S=D$ 时就是\stref{引理}{11.12}。另一个情形 $S=H$ 中，先对 $\nu:\partial H=\R\to\operatorname{Gr}(V)$ 作一个紧支的小变形，使之成为 $\nu(s)=\lambda(\psi(s))$，其中 $\lambda\in\Omega^-(\operatorname{Gr}(V),\allowbreak \Lambda_{\zeta,\allowbreak 0},\allowbreak \Lambda_{\zeta,\allowbreak 1})$，并对提升 $\nu^\#,\allowbreak \lambda^\#$ 作相应处理。这诱导 Fredholm 算子的变形，因此不改变\eref{11.26}左侧；也显然保持右侧，因为右侧与路径的任何选择无关。然后运用\stref{引理}{11.11}及 $i,\allowbreak o$ 的定义。还需要一个相关观察。取带出向无穷远点的上半平面，分别记为 $\bar H$ 与 $\bar\zeta$。赋予它反向路径 $\bar\nu^\#(s)=\nu^\#(-s)$，相应的 Lagrangian 子空间对仍与之前相同，$\Lambda_{\bar\zeta,\allowbreak k}=\Lambda_{\zeta,\allowbreak k}$。于是可以粘合 $\bar H\#H\cong D$，且与全部附加结构相容。应用\eref{11.9}得到
\begin{equation}\label{eq:11.27}
\begin{aligned}
n&=\operatorname{index}(D_{\bar H,\bar\nu})+i(\Lambda_{\bar\zeta,0}^\#,\Lambda_{\bar\zeta,1}^\#),\\
\lambda^{\mathrm{top}}(\bar\nu_z)&\cong\det(D_{\bar H,\bar\nu})\otimes o(\Lambda_{\bar\zeta,0}^\#,\Lambda_{\bar\zeta,1}^\#),
\end{aligned}
\end{equation}
其中 $z$ 为 $\partial\bar H$ 的任意点。这证明\eref{11.26}对 $\bar H$ 也成立。

一般策略如下。可以将紧化 $\hat S$ 视为 $k$ 个圆盘 $\hat D_1,\allowbreak \ldots,\allowbreak \hat D_k$ 与同亏格 $g$ 的闭曲面的连通和。以 $D_1,\allowbreak \ldots,\allowbreak D_k$ 表示这些圆盘的有限部分，即带边界标记点的部分。如\stref{引理}{11.7}证明的步骤 (i)，沿平行于 $\partial\hat D_j$ 的曲线捏合分离，得到
\[
\begin{aligned}
\operatorname{index}(D_{S,\nu})&=n(2-2g)+\sum_j(\operatorname{index}(D_{D_j,\nu_j})-2n),\\
\det(D_{S,\nu})&\cong\bigotimes_j\det(D_{D_j,\nu_j}),
\end{aligned}
\]
其中 $\nu_j$ 是 $\nu$ 在属于 $\hat C_j$ 的各分支上的限制。这就把问题归结为 $\hat S$ 为圆盘的情形，以下均作此假设。对每个出向无穷远点 $\zeta^+\in\Sigma^+$，选择 $\Omega^-(\operatorname{Gr}(V)^\#,\allowbreak \Lambda_{\zeta^+,\allowbreak 0}^\#,\allowbreak \Lambda_{\zeta^+,\allowbreak 1}^\#)$ 中的一条路径。取带相应边界数据的 $H$ 的副本，粘合到 $S$ 上，填补各出向端。若只有一个入向无穷远点，粘合后仍得到 $H$ 上的算子；由已证明的情形可见
\[
\begin{aligned}
i(\Lambda_{\zeta^-,0}^\#,\Lambda_{\zeta^-,1}^\#)&=\operatorname{index}(D_{S,\nu})+\sum_{\zeta^+}i(\Lambda_{\zeta^+,0}^\#,\Lambda_{\zeta^+,1}^\#),\\
o(\Lambda_{\zeta^-,0}^\#,\Lambda_{\zeta^-,1}^\#)&\cong\det(D_{S,\nu})\otimes\bigotimes_{\zeta^+}o(\Lambda_{\zeta^+,0}^\#,\Lambda_{\zeta^+,1}^\#).
\end{aligned}
\]
若没有入向无穷远点，或入向无穷远点至少有两个，则采取类似但稍复杂的办法，额外粘合 $\bar H$ 的副本，将
\SourcePage{174}
$S$ 完全紧化，再利用\eref{11.27}得出
\begin{equation}\label{eq:11.28}
\begin{aligned}
n={}&\sum_{\zeta^-}(n-i(\Lambda_{\zeta^-,0}^\#,\Lambda_{\zeta^-,1}^\#))
+\operatorname{index}(D_{S,\nu})+\sum_{\zeta^+}i(\Lambda_{\zeta^+,0}^\#,\Lambda_{\zeta^+,1}^\#),\\
\lambda^{\mathrm{top}}(\nu_{z_0})\cong{}&\bigotimes_{\zeta^-}\left(\lambda^{\mathrm{top}}(\nu_{z_j})\otimes o(\Lambda_{\zeta^-,0}^\#,\Lambda_{\zeta^-,1}^\#)^{-1}\right)\\
&\otimes\det(D_{S,\nu})\otimes\bigotimes_{\zeta^+}o(\Lambda_{\zeta^+,0}^\#,\Lambda_{\zeta^+,1}^\#).
\end{aligned}
\end{equation}
这里 $z_j$ 是 $S$ 的边界点。沿紧化曲面上的 Lagrangian 子空间环路作平行移动，可将各 $\lambda^{\mathrm{top}}(\nu_{z_j})$ 相互识别。由于该环路容许膜结构，它的 Maslov 指标为零，因而可定向，所以这里没有歧义。\znote{165}{这里用到了膜结构的全体数据：$\Lambda^\#$ 自带 Maslov 指标与 $\Z/2$ 障碍，故环路一旦提升到 $\operatorname{Gr}(V)^\#$ 中，它的 Maslov 指标与第二 Stiefel–Whitney 类便同时为零。前者使 $\lambda^{\mathrm{top}}$ 沿环路的平行移动没有单值性，后者使定向可以一致地搬运；两者合起来，才允许把各个 $\lambda^{\mathrm{top}}(\nu_{z_j})$ 相互等同。换言之，$\Z$ 的相位障碍与 $\Z/2$ 的符号障碍在此抵消。nLab 关于扭曲 spin 结构的条目说明了 $w_2$ 与 twist 如何合成新的障碍类。[来源：nLab "twisted spin structure"，https:/\allowbreak /\allowbreak ncatlab.org/\allowbreak nlab/\allowbreak show/\allowbreak twisted+\allowbreak spin+\allowbreak structure（2026-11-15 访问）]}这样就将\eref{11.28}化为与\eref{11.26}一致的形式。
\end{proof}

目前的表述有一个严重缺点：$\operatorname{Gr}(V)^\#$ 只在同伦等价意义下定义，使得抽象线性 Lagrangian 膜的概念带有一定的含混性。因此，下面用一个等价但更具几何意味的框架来代替，即分次与 Pin 结构。
